\documentclass[11pt,a4paper]{article}
\usepackage[utf8]{inputenc}
\usepackage[T1]{fontenc}
\usepackage[ngerman]{babel}
\usepackage{amsmath,amssymb,graphicx,booktabs,xcolor}
\usepackage[margin=25mm]{geometry}
\usepackage{tcolorbox,hyperref}

\definecolor{darkblue}{HTML}{16213e}
\definecolor{keygreen}{HTML}{2e7d32}

\title{\textcolor{darkblue}{Dok. 0017\\Diskant-Stimmzungen: Obertonmoden F3 bis C6}}
\author{Harmonikabau --- Technische Dokumentation}
\date{}

\begin{document}
\maketitle

\begin{tcolorbox}[colframe=red!60!black, colback=red!5!white]
\textbf{Hinweis:} Dieses Dokument bietet ein Erklärungsmodell.
Die interpolierte Mensur ist eine Schätzung --- reale Stimmplatten
variieren je nach Hersteller.
\end{tcolorbox}

\noindent\textit{\small Berechnungsskript: Die vollständige FEM-Berechnung mit allen Tabellen
ist im selben Verzeichnis auf GitHub abgelegt:
\textbf{pruefskript\_0017\_diskant.py} (W-Unabhängigkeit, gleichmäßige Mindestdicke,
kalibrierte Profilierung, $f_2$ vs.\ $f_H$, Vergleich Bass/Diskant).}

\noindent\textit{FEM-Berechnung der Biegemoden von Diskant-Stimmzungen mit variabler
Mensur (L = 37$\to$14\,mm, W = 3,7$\to$1,7\,mm). Kalibrierte Dicke für jeden
Halbton. Vergleich mit den Basszungen aus Dok.~0016.}

\section{Die Breite W hat keinen Einfluss auf die Tonhöhe}

Bei einem Euler-Bernoulli-Cantilever:
\begin{equation}
f = \frac{\beta^2}{2\pi L^2} \times t \times \sqrt{\frac{E}{12\rho}}
\end{equation}

Die Breite $W$ kürzt sich raus: $I = Wt^3/12$, $A = Wt$, also $I/A = t^2/12$.
Die Tonhöhe wird ausschließlich durch $t$ und $L$ bestimmt.
$W$ beeinflusst Masse ($\propto W$), Volumenstrom ($\propto W$),
Lautstärke und Torsionssteifigkeit ($\propto W^3$) --- aber nicht den Ton.

\begin{tcolorbox}[colframe=keygreen, colback=green!5!white]
\textbf{W bestimmt die Lautstärke und die Ansprache, nicht den Ton.}
Deshalb nehmen $W$ und $L$ stetig ab von tief zu hoch ---
die tiefen Töne brauchen mehr Volumenstrom.
\end{tcolorbox}

\section{Diskant-Mensur: L und W nehmen stetig ab, t wird profiliert}

Im Unterschied zum Bass variieren im Diskant Länge und Breite stetig.
Die Dicke am Fuß ($t_0$) ist immer etwas größer als die theoretische
Mindestdicke --- die Profilierung (Verdünnung zur Spitze) passt den Ton an.

\begin{center}
\begin{tabular}{lcccccccc}
\toprule
Ton & $f$ & $L$ & $W$ & $t_0$ & $t_\mathrm{tip}$ & Verd. & $f_2/f_1$ & $f_2$ \\
\midrule
F3 & 175 & 37,0 & 3,70 & 0,298 & 0,223 & 28\,\% & 5,32 & 348 \\
C3 & 131 & 29,0 & 3,05 & 0,133 & 0,104 & 24\,\% & 5,33 & 697 \\
C4 & 262 & 22,8 & 2,51 & 0,161 & 0,132 & 20\,\% & 5,45 & 1426 \\
C5 & 523 & 17,9 & 2,06 & 0,197 & 0,169 & 16\,\% & 5,63 & 2947 \\
C6 & 1047 & 14,0 & 1,70 & 0,244 & 0,217 & 12\,\% & 5,82 & 6094 \\
\bottomrule
\end{tabular}
\end{center}

\begin{tcolorbox}[colframe=keygreen, colback=green!5!white]
\textbf{$f_2/f_1 = 5{,}3$--$5{,}8$} (profiliert) --- niedriger als die 6,27
der gleichmäßigen Zunge, weil die Profilierung die Moden komprimiert.
Die leichte Profilierung im Diskant (12--28\,\%) senkt das Verhältnis
weniger als im Bass (30--50\,\%).
\end{tcolorbox}

\section{$f_2$ absolut: Wo trifft Mode~2 die Kammerresonanz?}

Im Bass liegt die kritische Zone bei A1--A2 ($f_2 = 345$--$689$\,Hz,
$f_H$ Bass $= 200$--$700$\,Hz). Im Diskant verschiebt sich alles nach oben:
$f_2$ liegt für C2--C3 bei 410--820\,Hz --- \textbf{unter} den Diskantkammern.
Erst ab C4 ($f_2 \approx 1640$\,Hz) beginnt $f_2$ den Diskantkammer-Bereich
zu treffen. Die kritische Zone im Diskant ist \textbf{C4--C5}
($f_2 = 1640$--$3280$\,Hz).

\section{Profilierung im Diskant: Weniger Spielraum}

Die Diskantzungen sind leicht profiliert: 12--25\,\% Verdünnung
(Basszungen: 30--50\,\%). Die dünnste Stelle liegt bei den hohen Tönen
an der Spitze, bei den tieferen eher in der Mitte. Die Profilierung senkt
$f_2/f_1$ um 0,4--0,9 (von 6,27 auf 5,4--5,8). Im Bass war die Senkung
1,0--2,3. Weniger Spielraum bei dünneren Zungen.

\section{Steifigkeit und Masse}

Die Steifigkeit steigt exponentiell: Faktor~100 von C2 (4,8\,N/m) bis
C6 (489\,N/m). Die Masse sinkt von 117\,mg auf 47\,mg. Hohe Töne:
steifer, leichter, höherer Schwellendruck, schnelleres Einschwingen.

\section{Vergleich Bass --- Diskant}

\textbf{Gemeinsam:} $f_2/f_1 = 6{,}27$ = konstant. Die Cantilever-Physik
ist identisch.

\textbf{Unterschied 1 --- Mensur:} Im Bass eine Zungenlänge für 2~Oktaven.
Im Diskant variieren $L$, $W$ und $t$ gleichzeitig.

\textbf{Unterschied 2 --- Profilierung:} Bass: 30--50\,\%, Effekt $-1{,}0$ bis
$-2{,}3$. Diskant: 12--25\,\%, Effekt $-0{,}4$ bis $-0{,}9$.

\textbf{Unterschied 3 --- Kritische Zone:} Bass: A1--A2 ($f_2$ im $f_H$-Bereich
der Basskammern). Diskant: C4--C5 ($f_2$ im $f_H$-Bereich der Diskantkammern).
Die unteren Diskant-Töne (F3--B3) sind unkritisch.

\section{Zusammenfassung}

\begin{enumerate}
\item \textbf{W hat keinen Einfluss auf die Tonhöhe.} $f \propto t/L^2$.
  $W$ bestimmt Lautstärke und Ansprache.
\item \textbf{$f_2/f_1 = 6{,}27$ (gleichmäßig) = identisch mit Bass.}
  Profiliert: 5,3--5,8 je nach Verdünnungstiefe.
\item \textbf{Reale Dicke am Fuß $>$ Mindestdicke.} Die Profilierung
  hebt $f_1$ an --- der Stimmer beginnt dicker und schleift herunter.
\item \textbf{Kritische Zone: C4--C5.} $f_2 = 1400$--$3000$\,Hz trifft
  die Diskantkammern. Untere Töne (F3--B3) unkritisch.
\item \textbf{Steifigkeit:} Faktor~100 über 4~Oktaven.
\end{enumerate}

\bigskip
\textit{Bass und Diskant gehorchen derselben Physik.
Die Mensur ändert die absoluten Frequenzen --- die Verhältnisse bleiben.}

\end{document}
